% Kommentare für den Editor (TexWorks/TexMakerX)
% !TeX encoding   = utf8
% !TeX spellcheck = de-DE

% --- LaTeX Vorlage ------------------------------------------------------
% speziell für einfache Dokumente wie Praktikumsprotokolle
%
% Autor: Matthias Pospiech (matthias@pospiech.eu)
% ------------------------------------------------------------------------

% Dokumentenklasse (Koma Script) -----------------------------------------
\documentclass[%
   %draft,     % Entwurfsstadium
   final,      % fertiges Dokument
   paper=a4, paper=portrait, pagesize=auto, % Papier Einstellungen
   fontsize=11pt, % Schriftgröße
   ngerman, % Sprache 
 ]{scrreprt} % Classes: scrartcl, scrreprt, scrbook

% ~~~~~~~~~~~~~~~~~~~~~~~~~~~~~~~~~~~~~~~~~~~~~~~~~~~~~~~~~~~~~~~~~~~~~~~~
% encoding
% ~~~~~~~~~~~~~~~~~~~~~~~~~~~~~~~~~~~~~~~~~~~~~~~~~~~~~~~~~~~~~~~~~~~~~~~~

% Encoding der Dateien (sonst funktionieren Umlaute nicht)
\usepackage[utf8]{inputenc}

% Encoding der Verzeichnisse (für Pfade mit Umlauten und Leerzeichne)
\usepackage[%
   extendedchars, encoding, multidot, space,
   %filenameencoding=latin1, % Windows XP, Vista, 7
   filenameencoding=utf8,   % Linux, OS X
]{grffile}

% ~~~~~~~~~~~~~~~~~~~~~~~~~~~~~~~~~~~~~~~~~~~~~~~~~~~~~~~~~~~~~~~~~~~~~~~~
% Pakete und Stile
% ~~~~~~~~~~~~~~~~~~~~~~~~~~~~~~~~~~~~~~~~~~~~~~~~~~~~~~~~~~~~~~~~~~~~~~~~
% Schriften
% ~~~~~~~~~~~~~~~~~~~~~~~~~~~~~~~~~~~~~~~~~~~~~~~~~~~~~~~~~~~~~~~~~~~~~~~~
% Fonts Fonts Fonts
% ~~~~~~~~~~~~~~~~~~~~~~~~~~~~~~~~~~~~~~~~~~~~~~~~~~~~~~~~~~~~~~~~~~~~~~~~

% immer laden:
\usepackage[T1]{fontenc} % T1 Schrift Encoding
\usepackage{textcomp}	 % Zusätzliche Symbole (Text Companion font extension)

% ~~~~~~~~~~~~~~~~~~~~~~~~~~~~~~~~~~~~~~~~~~~~~~~~~~~~~~~~~~~~~~~~~~~~~~~~
% Symbole
% ~~~~~~~~~~~~~~~~~~~~~~~~~~~~~~~~~~~~~~~~~~~~~~~~~~~~~~~~~~~~~~~~~~~~~~~~

\usepackage{amssymb}
\usepackage{mathcomp}


%% ==== Zusammengesetzte Schriften  (Sans + Serif) =======================

%% - Latin Modern
\usepackage{lmodern}
%% -------------------

%% - Bera Schriften
%\usepackage{bera}
%% -------------------

%% - Times, Helvetica, Courier (Word Standard...)
%\usepackage{mathptmx}
%\usepackage[scaled=.90]{helvet}
%\usepackage{courier}
%% -------------------

%% - Palantino , Helvetica, Courier
%\usepackage{mathpazo}
%\usepackage[scaled=.95]{helvet}
%\usepackage{courier}
%% -------------------

%% - Charter, Bera Sans
%\usepackage{charter}\linespread{1.05}
%\renewcommand{\sfdefault}{fvs}
%\usepackage[charter]{mathdesign}



%%%% =========== Typewriter =============

%\usepackage{courier}                   %% --- Courier
%\renewcommand{\ttdefault}{cmtl}        %% --- CmBright Typewriter Font
%\usepackage[%                          %% --- Luxi Mono (Typewriter)
%   scaled=0.9
%]{luximono}



% Pakete Laden
% ~~~~~~~~~~~~~~~~~~~~~~~~~~~~~~~~~~~~~~~~~~~~~~~~~~~~~~~~~~~~~~~~~~~~~~~~
% These packages must be loaded before all others
% (primarily because they are required by other packages)
% ~~~~~~~~~~~~~~~~~~~~~~~~~~~~~~~~~~~~~~~~~~~~~~~~~~~~~~~~~~~~~~~~~~~~~~~~
\usepackage{calc}
\usepackage{fixltx2e}	% Fix known LaTeX2e bugs

\usepackage[ngerman]{babel} 	% Sprache
\usepackage[dvipsnames, table]{xcolor} 	% Farben
\usepackage[ngerman, iso]{isodate} 	%Datumsformat

% ~~~~~~~~~~~~~~~~~~~~~~~~~~~~~~~~~~~~~~~~~~~~~~~~~~~~~~~~~~~~~~~~~~~~~~~~
% Bilder, Gleitumgebungen und Platzierung
% ~~~~~~~~~~~~~~~~~~~~~~~~~~~~~~~~~~~~~~~~~~~~~~~~~~~~~~~~~~~~~~~~~~~~~~~~

\usepackage[]{graphicx}					% Graphiken
\usepackage{epstopdf}		% konvertiert eps in pdf

% provides new floats and enables H float modifier option
\usepackage{float}
% Floats immer erst nach der Referenz setzen
\usepackage{flafter}
% Alel Floats werden vor der nächsten section ausgegeben
\usepackage[section]{placeins} 
%

\usepackage{tikz}
%tikz verwenden um Grafiken zu erstellen
\usepackage{pgfplots}
%pgfplots verwenden für Diagramme

% ~~~~~~~~~~~~~~~~~~~~~~~~~~~~~~~~~~~~~~~~~~~~~~~~~~~~~~~~~~~~~~~~~~~~~~~~
% Beschriftungen (captions)
% ~~~~~~~~~~~~~~~~~~~~~~~~~~~~~~~~~~~~~~~~~~~~~~~~~~~~~~~~~~~~~~~~~~~~~~~~

\usepackage{caption}
\usepackage{subcaption}

% ~~~~~~~~~~~~~~~~~~~~~~~~~~~~~~~~~~~~~~~~~~~~~~~~~~~~~~~~~~~~~~~~~~~~~~~~
% Math
% ~~~~~~~~~~~~~~~~~~~~~~~~~~~~~~~~~~~~~~~~~~~~~~~~~~~~~~~~~~~~~~~~~~~~~~~~

% Base Math Package
\usepackage[fleqn]{amsmath} 
% Warnt bei Benutzung von Befehlen die mit amsmath inkompatibel sind.
\usepackage[all, error]{onlyamsmath}
%Brüche schräg statt horizontal brechen
\usepackage{nicefrac}

% ~~~~~~~~~~~~~~~~~~~~~~~~~~~~~~~~~~~~~~~~~~~~~~~~~~~~~~~~~~~~~~~~~~~~~~~~
% Science
% ~~~~~~~~~~~~~~~~~~~~~~~~~~~~~~~~~~~~~~~~~~~~~~~~~~~~~~~~~~~~~~~~~~~~~~~~

% Einheiten und Zahlenformatierung
\usepackage{siunitx}

% ~~~~~~~~~~~~~~~~~~~~~~~~~~~~~~~~~~~~~~~~~~~~~~~~~~~~~~~~~~~~~~~~~~~~~~~~
% Tables (Tabular)
% ~~~~~~~~~~~~~~~~~~~~~~~~~~~~~~~~~~~~~~~~~~~~~~~~~~~~~~~~~~~~~~~~~~~~~~~~

\usepackage{booktabs}
\usepackage{ltxtable} % Longtable + tabularx

% ~~~~~~~~~~~~~~~~~~~~~~~~~~~~~~~~~~~~~~~~~~~~~~~~~~~~~~~~~~~~~~~~~~~~~~~~
% text related packages
% ~~~~~~~~~~~~~~~~~~~~~~~~~~~~~~~~~~~~~~~~~~~~~~~~~~~~~~~~~~~~~~~~~~~~~~~~

\usepackage{url}            % Befehl \url{...}
\usepackage{enumitem}		% Kompakte Listen

% Neue Befehle: \Centering, \RaggedLeft, and \RaggedRight, ... 
\usepackage{ragged2e}


% ~~~~~~~~~~~~~~~~~~~~~~~~~~~~~~~~~~~~~~~~~~~~~~~~~~~~~~~~~~~~~~~~~~~~~~~~
% Citations
% ~~~~~~~~~~~~~~~~~~~~~~~~~~~~~~~~~~~~~~~~~~~~~~~~~~~~~~~~~~~~~~~~~~~~~~~~

%\usepackage[
%	style=alphabetic, % Loads the bibliography and the citation style 
%	natbib=true, % define natbib compatible cite commands
%]{biblatex}	
% Other options:
%	style=numeric, % 
%	style=numeric-comp,    % [1–3, 7, 8]
%	style=numeric-verb,    % [2]; [5]; [6]

\usepackage{csquotes}

% ~~~~~~~~~~~~~~~~~~~~~~~~~~~~~~~~~~~~~~~~~~~~~~~~~~~~~~~~~~~~~~~~~~~~~~~~
% layout packages
% ~~~~~~~~~~~~~~~~~~~~~~~~~~~~~~~~~~~~~~~~~~~~~~~~~~~~~~~~~~~~~~~~~~~~~~~~
%
% Befehle für 1,5 und 2 zeilig: 
% \singlespacing, \onehalfspacing und \doublespacing
\usepackage{setspace}

% ~~~~~~~~~~~~~~~~~~~~~~~~~~~~~~~~~~~~~~~~~~~~~~~~~~~~~~~~~~~~~~~~~~~~~~~~
% Kopf und Fusszeile
% ~~~~~~~~~~~~~~~~~~~~~~~~~~~~~~~~~~~~~~~~~~~~~~~~~~~~~~~~~~~~~~~~~~~~~~~~

% Kopf und Fusszeile mit scrpage2 einstellen
\usepackage[automark, komastyle, nouppercase]{scrpage2}

% ~~~~~~~~~~~~~~~~~~~~~~~~~~~~~~~~~~~~~~~~~~~~~~~~~~~~~~~~~~~~~~~~~~~~~~~~
% pdf packages
% ~~~~~~~~~~~~~~~~~~~~~~~~~~~~~~~~~~~~~~~~~~~~~~~~~~~~~~~~~~~~~~~~~~~~~~~~

% Include pages from external PDF documents in LaTeX documents
\usepackage{pdfpages} 

% Optischer Randausgleich mit pdfTeX
\usepackage{microtype}

% Einstellungen und Layoutstile 
% ~~~~~~~~~~~~~~~~~~~~~~~~~~~~~~~~~~~~~~~~~~~~~~~~~~~~~~~~~~~~~~~~~~~~~~~~
% Colors
% ~~~~~~~~~~~~~~~~~~~~~~~~~~~~~~~~~~~~~~~~~~~~~~~~~~~~~~~~~~~~~~~~~~~~~~~~
\definecolor{sectioncolor}{RGB}{0, 0, 0}     % black

% ~~~~~~~~~~~~~~~~~~~~~~~~~~~~~~~~~~~~~~~~~~~~~~~~~~~~~~~~~~~~~~~~~~~~~~~~
% text related 
% ~~~~~~~~~~~~~~~~~~~~~~~~~~~~~~~~~~~~~~~~~~~~~~~~~~~~~~~~~~~~~~~~~~~~~~~~

%% style of URL
\urlstyle{tt}


% Keine hochgestellten Ziffern in der Fussnote (KOMA-Script-spezifisch):
\deffootnote{1.5em}{1em}{\makebox[1.5em][l]{\thefootnotemark}}

% Limit space of footnotes to 10 lines
\setlength{\dimen\footins}{10\baselineskip}

% prevent continuation of footnotes 
% at facing page
\interfootnotelinepenalty=10000 

% ~~~~~~~~~~~~~~~~~~~~~~~~~~~~~~~~~~~~~~~~~~~~~~~~~~~~~~~~~~~~~~~~~~~~~~~~
% Science
% ~~~~~~~~~~~~~~~~~~~~~~~~~~~~~~~~~~~~~~~~~~~~~~~~~~~~~~~~~~~~~~~~~~~~~~~~

\sisetup{%
	mode = math, detect-family, detect-weight,	
	exponent-product = \cdot,
	number-unit-separator=\text{\,},
	output-decimal-marker={,},
}

% ~~~~~~~~~~~~~~~~~~~~~~~~~~~~~~~~~~~~~~~~~~~~~~~~~~~~~~~~~~~~~~~~~~~~~~~~
% Citations / Style of Bibliography
% ~~~~~~~~~~~~~~~~~~~~~~~~~~~~~~~~~~~~~~~~~~~~~~~~~~~~~~~~~~~~~~~~~~~~~~~~

% Kommentar entfernene wenn biblatex geladen wird
% \IfPackageLoaded{biblatex}{%
	\ExecuteBibliographyOptions{%
%--- Backend --- --- ---
	backend=bibtex,  % (bibtex, bibtex8, biber)
	bibwarn=true, %
	bibencoding=ascii, % (ascii, inputenc, <encoding>)
%--- Sorting --- --- ---
	sorting=nty, % Sort by name, title, year.
	% other options: 
	% nty        Sort by name, title, year.
	% nyt        Sort by name, year, title.
	% nyvt       Sort by name, year, volume, title.
	% anyt       Sort by alphabetic label, name, year, title.
	% anyvt      Sort by alphabetic label, name, year, volume, title.
	% ynt        Sort by year, name, title.
	% ydnt       Sort by year (descending), name, title.
	% none       Do not sort at all. All entries are processed in citation order.
	% debug      Sort by entry key. This is intended for debugging only.
	%
	sortcase=true,
	sortlos=los, % (bib, los) The sorting order of the list of shorthands
	sortcites=false, % do/do not sort citations according to bib	
%--- Dates --- --- ---
	date=comp,  % (short, long, terse, comp, iso8601)
%	origdate=
%	eventdate=
%	urldate=
%	alldates=
	datezeros=true, %
	dateabbrev=true, %
%--- General Options --- --- ---
	maxnames=1,
	minnames=1,
%	maxbibnames=99,
%	maxcitenames=1,
%	autocite= % (plain, inline, footnote, superscript) 
	autopunct=true,
	language=auto,
	babel=none, % (none, hyphen, other, other*)
	block=none, % (none, space, par, nbpar, ragged)
	notetype=foot+end, % (foot+end, footonly, endonly)
	hyperref=true, % (true, false, auto)
	backref=true,
	backrefstyle=three, % (none, three, two, two+, three+, all+)
	backrefsetstyle=setonly, %
	indexing=false, % 
	% options:
	% true       Enable indexing globally.
	% false      Disable indexing globally.
	% cite       Enable indexing in citations only.
	% bib        Enable indexing in the bibliography only.
	refsection=none, % (part, chapter, section, subsection)
	refsegment=none, % (none, part, chapter, section, subsection)
	abbreviate=true, % (true, false)
	defernumbers=false, % 
	punctfont=false, % 
	arxiv=abs, % (ps, pdf, format)	
%--- Style Options --- --- ---	
% The following options are provided by the standard styles
	isbn=false,%
	url=false,%
	doi=false,%
	eprint=false,%	
	}%	
	
	% change alpha label to be without +	
	\renewcommand*{\labelalphaothers}{}
	
	% change 'In: <magazine>" to "<magazine>"
	\renewcommand*{\intitlepunct}{}
	\DefineBibliographyStrings{german}{in={}}
	
	% make names capitalized \textsc{}
	\renewcommand{\mkbibnamefirst}{\textsc}
	\renewcommand{\mkbibnamelast}{\textsc}
	
	% make volume and number look like 
	% 'Bd. 33(14): '
	\renewbibmacro*{volume+number+eid}{%
	  \setunit{\addcomma\space}%
	  \bibstring{volume}% 
	  \setunit{\addspace}%
	  \printfield{volume}%
	  \iffieldundef{number}{}{% 
	    \printtext[parens]{%
	      \printfield{number}%
	    }%
	  }%
	  \setunit{\addcomma\space}%
	  \printfield{eid}
	  %\setunit{\addcolon\space}%
	  }	

	% <authors>: <title>
	\renewcommand*{\labelnamepunct}{\addcolon\space}
	% make ': ' before pages
	\renewcommand*{\bibpagespunct}{\addcolon\space}
	% names delimiter ';' instead of ','
	%\renewcommand*{\multinamedelim}{\addsemicolon\space}

	% move date before issue
	\renewbibmacro*{journal+issuetitle}{%
	  \usebibmacro{journal}%
	  \setunit*{\addspace}%
	  \iffieldundef{series}
	    {}
	    {\newunit
	     \printfield{series}%
	     \setunit{\addspace}}%
	  %
	  \usebibmacro{issue+date}%
	  \setunit{\addcolon\space}%
	  \usebibmacro{issue}%
	  \setunit{\addspace}%
	  \usebibmacro{volume+number+eid}%
	  \newunit}

	% print all names, even if maxnames = 1
	\DeclareCiteCommand{\citeauthors}
	  {
	   \defcounter{maxnames}{1000}
	   \boolfalse{citetracker}%
	   \boolfalse{pagetracker}%
	   \usebibmacro{prenote}}
	  {\ifciteindex
	     {\indexnames{labelname}}
	     {}%
	   \printnames{labelname}}
	  {\multicitedelim}
	  {\usebibmacro{postnote}}

}%

% ~~~~~~~~~~~~~~~~~~~~~~~~~~~~~~~~~~~~~~~~~~~~~~~~~~~~~~~~~~~~~~~~~~~~~~~~
% figures, placement, floats and captions
% ~~~~~~~~~~~~~~~~~~~~~~~~~~~~~~~~~~~~~~~~~~~~~~~~~~~~~~~~~~~~~~~~~~~~~~~~

% Make float placement easier
\renewcommand{\floatpagefraction}{.75} % vorher: .5
\renewcommand{\textfraction}{.1}       % vorher: .2
\renewcommand{\topfraction}{.8}        % vorher: .7
\renewcommand{\bottomfraction}{.5}     % vorher: .3
\setcounter{topnumber}{3}        % vorher: 2
\setcounter{bottomnumber}{2}     % vorher: 1
\setcounter{totalnumber}{5}      % vorher: 3

%% ~~~ Captions ~~~~~~~~~~~~~~~~~~~~~~~~~~~~~~~~~~~~~~~~~~~~~~~~~~~~~~~~~~
% Style of captions
\DeclareCaptionStyle{captionStyleTemplateDefault}
[ % single line captions
   justification = centering
]
{ % multiline captions
% -- Formatting
   format      = plain,  % plain, hang
   indention   = 0em,    % indention of text 
   labelformat = default,% default, empty, simple, brace, parens
   labelsep    = colon,  % none, colon, period, space, quad, newline, endash
   textformat  = simple, % simple, period
% -- Justification
   justification = justified, %RaggedRight, justified, centering
   singlelinecheck = true, % false (true=ignore justification setting in single line)
% -- Fonts
   labelfont   = {small,bf},
   textfont    = {small,rm},
% valid values:
% scriptsize, footnotesize, small, normalsize, large, Large
% normalfont, ip, it, sl, sc, md, bf, rm, sf, tt
% singlespacing, onehalfspacing, doublespacing
% normalcolor, color=<...>
%
% -- Margins and further paragraph options
   margin = 10pt, %.1\textwidth,
   % width=.8\linewidth,
% -- Skips
   skip     = 10pt, % vertical space between the caption and the figure
   position = auto, % top, auto, bottom
% -- Lists
   % list=no, % suppress any entry to list of figure 
   listformat = subsimple, % empty, simple, parens, subsimple, subparens
% -- Names & Numbering
   % figurename = Abb. %
   % tablename  = Tab. %
   % listfigurename=
   % listtablename=
   % figurewithin=chapter
   % tablewithin=chapter
%-- hyperref related options
	hypcap=true, % (true, false) 
	% true=all hyperlink anchors are placed at the 
	% beginning of the (floating) environment
	%
	hypcapspace=0.5\baselineskip
}

% apply caption style
\captionsetup{
	style = captionStyleTemplateDefault % base
}

% Predefinded skip setup for different floats
\captionsetup[table]{position=top}
\captionsetup[figure]{position=bottom}


% options for subcaptions
\captionsetup[sub]{ %
	style = captionStyleTemplateDefault, % base
	skip=6pt,
	margin=5pt,
	labelformat = parens,% default, empty, simple, brace
	labelsep    = space,
	list=false,
	hypcap=false
}

% ~~~~~~~~~~~~~~~~~~~~~~~~~~~~~~~~~~~~~~~~~~~~~~~~~~~~~~~~~~~~~~~~~~~~~~~~
% layout 
% ~~~~~~~~~~~~~~~~~~~~~~~~~~~~~~~~~~~~~~~~~~~~~~~~~~~~~~~~~~~~~~~~~~~~~~~~


%% Paragraph Separation =================================
\KOMAoptions{%
   parskip=absolute, % do not change indentation according to fontsize
   parskip=false     % indentation of 1em
   % parskip=half    % parksip of 1/2 line 
}%

%% line spacing =========================================
%\onehalfspacing	% 1,5-facher Abstand
%\doublespacing		% 2-facher Abstand

%% page layout ==========================================

\raggedbottom     % Variable Seitenhoehen zulassen

% Koma Script text area layout
\KOMAoptions{%
   DIV=11,% (Size of Text Body, higher values = greater textbody)
   BCOR=5mm% (Bindekorrektur)
}%

%%% === Page Layout  Options ===
\KOMAoptions{% (most options are for package typearea)
   % twoside=true, % two side layout (alternating margins, standard in books)
   twoside=false, % single side layout 
   %
   headlines=2.1,%
}%

%\KOMAoptions{%
%      headings=noappendixprefix % chapter in appendix as in body text
%      ,headings=nochapterprefix  % no prefix at chapters
%      % ,headings=appendixprefix   % inverse of 'noappendixprefix'
%      % ,headings=chapterprefix    % inverse of 'nochapterprefix'
%      % ,headings=openany   % Chapters start at any side
%      % ,headings=openleft  % Chapters start at left side
%      ,headings=openright % Chapters start at right side      
%}%


% reloading of typearea, necessary if setting of spacing changed
\typearea[current]{last}

%% Zitate
\renewcommand*{\dictumwidth}{0.5\textwidth}

% ~~~~~~~~~~~~~~~~~~~~~~~~~~~~~~~~~~~~~~~~~~~~~~~~~~~~~~~~~~~~~~~~~~~~~~~~
% Titlepage
% ~~~~~~~~~~~~~~~~~~~~~~~~~~~~~~~~~~~~~~~~~~~~~~~~~~~~~~~~~~~~~~~~~~~~~~~~
\KOMAoptions{%
   % titlepage=true %
   titlepage=false %
}%

% ~~~~~~~~~~~~~~~~~~~~~~~~~~~~~~~~~~~~~~~~~~~~~~~~~~~~~~~~~~~~~~~~~~~~~~~~
% head and foot lines
% ~~~~~~~~~~~~~~~~~~~~~~~~~~~~~~~~~~~~~~~~~~~~~~~~~~~~~~~~~~~~~~~~~~~~~~~~

% \pagestyle{scrheadings} % Seite mit Headern
\pagestyle{scrplain} % Seiten ohne Header

% loescht voreingestellte Stile
\clearscrheadings
\clearscrplain
%
% Was steht wo...
% Bei headings:
%   % Oben aussen: Kapitel und Section
%   % Unten aussen: Seitenzahl
%   \ohead{\pagemark}
%   \ihead{\headmark}
%   \ofoot[\pagemark]{} % Außen unten: Seitenzahlen bei plain
% Bei Plain:
\cfoot[\pagemark]{\pagemark} % Mitte unten: Seitenzahlen bei plain


% Angezeigte Abschnitte im Header
% \automark[section]{chapter} %[rechts]{links}
\automark[subsection]{section} %[rechts]{links}

% ~~~~~~~~~~~~~~~~~~~~~~~~~~~~~~~~~~~~~~~~~~~~~~~~~~~~~~~~~~~~~~~~~~~~~~~~
% headings / page opening
% ~~~~~~~~~~~~~~~~~~~~~~~~~~~~~~~~~~~~~~~~~~~~~~~~~~~~~~~~~~~~~~~~~~~~~~~~
\setcounter{secnumdepth}{3}

\KOMAoptions{%
%%%% headings
   % headings=small  % Small Font Size, thin spacing above and below
   % headings=normal % Medium Font Size, medium spacing above and below
   headings=big % Big Font Size, large spacing above and below
}%

% Titelzeile linksbuendig, haengend
\renewcommand*{\raggedsection}{\raggedright} 

% ~~~~~~~~~~~~~~~~~~~~~~~~~~~~~~~~~~~~~~~~~~~~~~~~~~~~~~~~~~~~~~~~~~~~~~~~
% fonts of headings
% ~~~~~~~~~~~~~~~~~~~~~~~~~~~~~~~~~~~~~~~~~~~~~~~~~~~~~~~~~~~~~~~~~~~~~~~~
\setkomafont{sectioning}{\normalfont\sffamily} % \rmfamily
\setkomafont{descriptionlabel}{\itshape}
\setkomafont{pageheadfoot}{\normalfont\normalcolor\small\sffamily}
\setkomafont{pagenumber}{\normalfont\sffamily}

%%% --- Titlepage ---
%\setkomafont{subject}{}
%\setkomafont{subtitle}{}
%\setkomafont{title}{}

% ~~~~~~~~~~~~~~~~~~~~~~~~~~~~~~~~~~~~~~~~~~~~~~~~~~~~~~~~~~~~~~~~~~~~~~~~
% settings and layout of TOC, LOF, 
% ~~~~~~~~~~~~~~~~~~~~~~~~~~~~~~~~~~~~~~~~~~~~~~~~~~~~~~~~~~~~~~~~~~~~~~~~
\setcounter{tocdepth}{3} % Depth of TOC Display

% ~~~~~~~~~~~~~~~~~~~~~~~~~~~~~~~~~~~~~~~~~~~~~~~~~~~~~~~~~~~~~~~~~~~~~~~~
% Tabellen
% ~~~~~~~~~~~~~~~~~~~~~~~~~~~~~~~~~~~~~~~~~~~~~~~~~~~~~~~~~~~~~~~~~~~~~~~~

%%% -| Neue Spaltendefinitionen 'columntypes' |--
%
% Belegte Spaltentypen:
% l - links
% c - zentriert
% r - rechts
% p,m,b  - oben, mittig, unten
% X - tabularx Auto-Spalte

% um Tabellenspalten mit Flattersatz zu setzen, muss \\ vor
% (z.B.) \raggedright geschuetzt werden:
\newcommand{\PreserveBackslash}[1]{\let\temp=\\#1\let\\=\temp}

% Spalten mit Flattersatz und definierte Breite:
% m{} -> mittig
% p{} -> oben
% b{} -> unten
%
% Linksbuendig:
\newcolumntype{v}[1]{>{\PreserveBackslash\RaggedRight\hspace{0pt}}p{#1}}
\newcolumntype{M}[1]{>{\PreserveBackslash\RaggedRight\hspace{0pt}}m{#1}}
% % Rechtsbuendig :
% \newcolumntype{R}[1]{>{\PreserveBackslash\RaggedLeft\hspace{0pt}}m{#1}}
% \newcolumntype{S}[1]{>{\PreserveBackslash\RaggedLeft\hspace{0pt}}p{#1}}
% % Zentriert :
% \newcolumntype{Z}[1]{>{\PreserveBackslash\Centering\hspace{0pt}}m{#1}}
% \newcolumntype{A}[1]{>{\PreserveBackslash\Centering\hspace{0pt}}p{#1}}

\newcolumntype{Y}{>{\PreserveBackslash\RaggedLeft\hspace{0pt}}X}

%-- Einstellungen für Tabellen ----------
\providecommand\tablestyle{%
  \renewcommand{\arraystretch}{1.4} % Groessere Abstaende zwischen Zeilen
  \normalfont\normalsize            %
  \sffamily\small           % Serifenlose und kleine Schrift
  \centering%                       % Tabelle zentrieren
}

%--Einstellungen für Tabellen ----------

\colorlet{tablesubheadcolor}{gray!40}
\colorlet{tableheadcolor}{gray!25}
\colorlet{tableblackheadcolor}{black!60}
\colorlet{tablerowcolor}{gray!15.0}


% ~~~~~~~~~~~~~~~~~~~~~~~~~~~~~~~~~~~~~~~~~~~~~~~~~~~~~~~~~~~~~~~~~~~~~~~~
% pdf packages
% ~~~~~~~~~~~~~~~~~~~~~~~~~~~~~~~~~~~~~~~~~~~~~~~~~~~~~~~~~~~~~~~~~~~~~~~~

% ~~~~~~~~~~~~~~~~~~~~~~~~~~~~~~~~~~~~~~~~~~~~~~~~~~~~~~~~~~~~~~~~~~~~~~~~
% fix remaining problems
% ~~~~~~~~~~~~~~~~~~~~~~~~~~~~~~~~~~~~~~~~~~~~~~~~~~~~~~~~~~~~~~~~~~~~~~~~




% ~~~~~~~~~~~~~~~~~~~~~~~~~~~~~~~~~~~~~~~~~~~~~~~~~~~~~~~~~~~~~~~~~~~~~~~~
% Eigene Befehle
% ~~~~~~~~~~~~~~~~~~~~~~~~~~~~~~~~~~~~~~~~~~~~~~~~~~~~~~~~~~~~~~~~~~~~~~~~
% -- new commands --
\providecommand{\abs}[1]{\lvert#1\rvert}
\providecommand{\Abs}[1]{\left\lvert#1\right\rvert}
\providecommand{\norm}[1]{\left\Vert#1\right\Vert}
\providecommand{\Trace}[1]{\ensuremath{\Tr\{\,#1\,\}}} % Trace /Spur
%

\renewcommand{\d}{\partial\mspace{2mu}} % partial diff
\newcommand{\td}{\,\mathrm{d}}	% total diff

\newcommand{\Ham}{\mathcal{H}}    % Hamilton
\newcommand{\Prob}{\mathscr{P}}    % Hamilton
\newcommand{\unity}{\mathds{1}}   % Real

\renewcommand{\i}{\mathrm{i}}   % imagin�re Einheit



% -- New Operators --
\DeclareMathOperator{\rot}{rot}
\DeclareMathOperator{\grad}{grad}
\DeclareMathOperator{\Tr}{Tr}
\DeclareMathOperator{\const}{const}
\DeclareMathOperator{\e}{e} 			% exponatial Function


%\renewcommand{\labelenumi}{\alph{enumi})}
%\renewcommand{\labelenumii}{\roman{enumii})}
\setlength{\parindent}{0pt}
% ~~~~~~~~~~~~~~~~~~~~~~~~~~~~~~~~~~~~~~~~~~~~~~~~~~~~~~~~~~~~~~~~~~~~~~~~
% Eigene Befehle
% ~~~~~~~~~~~~~~~~~~~~~~~~~~~~~~~~~~~~~~~~~~~~~~~~~~~~~~~~~~~~~~~~~~~~~~~~
% Silbentrennung hinzufügen als 
% Sil-ben-tren-nung 
\hyphenation{}

\listfiles % schreibt alle verwendeten Dateien in die log Datei

%% Dokument Beginn %%%%%%%%%%%%%%%%%%%%%%%%%%%%%%%%%%%%%%%%%%%%%%%%%%%%%%%%
\begin{document}
%\includepdf{content/Deckblatt_WI}
% Automatische Titelseite

\subject{Physik II}
\title{Formelsammlung\\}
\author{\includegraphics[width=6.5cm]{images/logo}\\ \\  David Mändlen \& Ulf Schmelzer}
\date{\today}
\maketitle

% Manuelle Titelseite

%\begin{titlepage}
%   \mbox{}\vspace{5\baselineskip}\\
%   \sffamily\huge
%   \centering
%   % Titel
%   \LaTeXe{}-Vorlage von Matthias Pospiech
%   \vspace{3\baselineskip}\\
%   \rmfamily\Large
%   Leibniz Universität Hannover
%   \vspace{2\baselineskip}\\
%   \rmfamily\Large
%   Matthias Pospiech
%   \vspace{1\baselineskip}\\
%   \today
%\end{titlepage}


%\clearpage

\null

\vfill

\dictum[Dr. Axel Stoll]{Magie ist Physik durch Wollen!}

\clearpage

\tableofcontents

% Testdokumente (auskommentieren!)
%\input{content/demo/demo.tex}
%\input{content/demo/latexexample.tex}

% in diese Datei gehört der Inhalt des Dokumentes:
%\input{content/Fehlerrechnung.tex}
%\input{content/Vorbereitung.tex}
\chapter{Wellenlehre}
	\section{Grundlagen}
		\subsection{Ausbreitungsgeschwindigkeit der Welle $c$}
			\begin{align*}
			c &= \frac{\lambda}{T} = \lambda f = \frac{\Delta l}{\Delta t} 
			& \frac{c}{c_0}&=\sqrt{\frac{T}{T_0}}
			\end{align*}

			\begin{table}[h]
			\begin{tabular}{lll}
			$\lambda\dots$ Wellenlänge & $\Delta l\dots$ Länge & $\Delta t\dots$ Zeit\\
			$T_0\dots$ Vergleichstemperatur & $c_0\dots$ Vergleichsgeschwindigkeit &\\
			\end{tabular}
			\end{table}
		
		\subsection{Frequenzdifferenz}
			\[\frac{\Delta f}{f}=\frac{\Delta \lambda}{\lambda}\]
        
		\subsection{Zeitverzögerung}
			\[ \Delta t = \frac{x}{c} \]
			
		\subsection{Grundgleichungen}
			\begin{align*}
				y(x,t) & = \hat{y} \cos [\omega(t-\frac{x}{c})] &oder& & y(x,t) &= \hat{y} \cos(\omega t-kx) 
			\end{align*}
			wobei - Verschiebung nach rechts (pos x Richtung) \newline
			und + Verschiebung nach links
        
		\subsection{Wellenzahl $k$}
			\[ k  = \frac{\omega}{c} = \frac{2\pi f}{c} = \frac{2\pi}{\lambda} \]
			Infos: \newline
			Lichtwellen von $\lambda_{min} = 380 \mathrm{nm}$ (Ultra Violett) bis $\lambda_{max} = 780\mathrm{nm}$ rot \newline
			Lichtfrequenzen von $f_{max} = 790\mathrm{THz}$ bis $f_{min} = 390 \mathrm{THz}$
        
		\subsection{En.Stromdichte/Intensität I / S}
			oder Energiestromdichte
			\begin{align*}
				I &= c \omega	& I &= \frac12 c \varrho y^2 \omega^2 
			\end{align*} 

			\begin{table}[h]
			\begin{tabular}{ll}
			$\omega \dots$ Energiedichte & $c \dots$ Wellenbewegung\\
			\end{tabular}
			\end{table}

		\subsection{Grund- und Oberschwingungen}
			\[
				f_n=(n+1)\cdot\frac{c}{2l}
			\]
			\begin{table}[h]
			\begin{tabular}{lll}
				$f_n\dots$ Schwingungsfrequenz & $c\dots$ Phasengeschwindigkeit & $l\dots$ Saitenlänge\\
			\end{tabular}
			\end{table}
		
		\subsection{Wellenlänge im Medium}
			\[
				\lambda_{\mathrm{m}}=\frac{c_0}{\sqrt{\varepsilon_{\mathrm{r}}f}}
			\]
			
			\begin{table}[h]
			\begin{tabular}{lll}
				$\lambda_{\mathrm{m}}\dots$ Wellenlänge im Medium 
				& $c_0\dots$ Ausbreitungsgeschwindigkeit 
				& $\varepsilon_{\mathrm{r}}\dots$ Permittivitätszahl\\
			\end{tabular}
			\end{table}

		\subsection{Antennen}
			\begin{table}[h]
			\begin{tabular}{ll}
				Antennenart & Formel\\
				Dipol & $l=\frac{\lambda}{2}=\frac{c}{2f}$\\
				Stab & $l=\frac{\lambda}{4}=\frac{c}{4f}$\\
			\end{tabular}
			\end{table}
			\begin{table}[h]
			\begin{tabular}{lll}
				$c\dots$ Lichtgeschwindigkeit & $l\dots$ Antennenlänge\\
			\end{tabular}
			\end{table}
	\section{Wellenüberlagerung}
		\subsection{Grundlagen}
			\begin{align*}
				y_1&=\hat{y}\cos(\omega_1t-k_1x) & \omega_1\approx\omega_2\Rightarrow\omega&=\frac{\omega_1+\omega_2}{2}\\
				y_2&=\hat{y}\cos(\omega_2t-k_2x) & k_1\approx k_2\Rightarrow k&=\frac{k_1+k_2}{2}\\
				y=y_1+y_2&=2\hat{y}\cos(\omega t-kx)\cos\left(\frac{\Delta\omega}{2}t-\frac{\Delta k}{2}x\right)
			\end{align*}

		\subsubsection{Phasen-/Gruppengeschwindigkeit}
			\begin{align*}
				c&=\frac{\omega}{k} 
				& c_{\mathrm{Gr}}&=\frac{\mathrm{d}x}{\mathrm{d}t}=\frac{\Delta\omega}{\Delta k}
				=\frac{\mathrm{d}\omega}{\mathrm{d}t}=c-\lambda\cdot\frac{\mathrm{d}c}{\mathrm{d}\lambda}
				=\frac{f_1-f_2}{\frac{1}{\lambda_1}-\frac{1}{\lambda_2}}
			\end{align*}

		\subsection{Dispersion}
			\subsubsection{Fallunterscheidung}
				\begin{table}[h]
				\begin{tabular}{ll}
					Kriterium & Dispersionsart\\
					\toprule
					$c_{\mathrm{Gr}} = c$ & Keine Dispersion\\
					$c_{\mathrm{Gr}} < c$ & Normale Dispersion\\
					$c_{\mathrm{Gr}} > c$ & Anomale Dispersion\\
				\end{tabular}
				\end{table}

			\subsubsection{Biegeschwingung}
				\begin{align*}
					c&=\lambda_0\cdot f_0=\sqrt{\frac{\pi df_0}{2}}\cdot\sqrt[4]{\frac{E}{\varrho}}
					& \lambda_0&=2l
					& f_0=\frac{c}{2l}
				\end{align*}

				\begin{table}[h]
				\begin{tabular}{lll}
					$\lambda_0\dots$ Wellenlänge Grundschwingung & $f_0\dots$ Frequenz Grundschwingung & $d\dots$ Durchmesser Stab\\
					$E\dots$ Elastizitätsmodul,$\left[\frac{\mathrm{N}}{\mathrm{m}^2}\right]$ & $\varrho\dots$ Werkstoffdichte\\
				\end{tabular}
				\end{table}

			\subsubsection{Interferometer (Michelson)}
				\begin{align*}
					\lambda&=2\Delta s 
					& t_{\mathrm{Luft}}&=\frac{x}{c_0} 
					& t_{\mathrm{Platte}}&=\frac{x-l-\Delta x}{c_0}+\frac{ln}{c_0}\\
					x&=x-l-\Delta x+ln\Rightarrow & n&=1+\frac{\Delta x}{l} 
					& c&=\frac{c_0}{n}
				\end{align*}

				\begin{table}[h]
				\begin{tabular}{lll}
				$t\dots$ Laufzeit im Medium & $x\dots$ Entfernung Teiler zu Spiegel & $\Delta x$ Spiegelverschiebung $\mathrm{S}_{\mathrm{rechts}}$\\
				$l\dots$ Probendicke & $c_0\dots$ Lichtgeschwindigkeit in Luft & $n$ Brechungsindex\\
				$c\dots$ Lichtgeschw. (Medium) & $\Delta s$ Spiegelverschiebung $\mathrm{S}_{\mathrm{oben}}$\\
				\end{tabular}
				\end{table}
			
			\subsubsection{Maxwell-Relation}
				\begin{align*}
					n&=\frac{c_0}{c} & n&=\sqrt{\varepsilon_{\mathrm{r}}}
				\end{align*}

				\begin{table}[h]
				\begin{tabular}{lll}
					$n\dots$ Brechnungsindex & $c\dots$ Ausbreitungsgeschw. (Medium) & $c_0\dots$ Lichtgeschw.\\
				\end{tabular}
				\end{table}

		\subsection{Leiterarten}
			\subsubsection{Hohlleiter}
				\begin{align*}
					\omega(k)&=c_0\sqrt{k^2+\left(\frac{\pi}{a}\right)^2}
					& c&=\frac{c_0}{\sqrt{1-\left(\frac{c_0}{2af}\right)^2}}
					& c_{\mathrm{Gr}}&=\frac{\mathrm{d}\omega}{\mathrm{d}t}=c_0\sqrt{1-\left(\frac{c_0}{2af}\right)^2}
				\end{align*}

			\subsubsection{Quarzglas}
				\begin{align*}
					n&=\frac{c_0}{c}
					& c_{\mathrm{Gr}}&=c-\lambda\frac{\mathrm{d}c}{\mathrm{d}\lambda}=\frac{c_0}{n_{\mathrm{Gr}}}
					& n_{\mathrm{Gr}}&=n-\lambda\frac{\mathrm{d}n}{\mathrm{d}\lambda}
				\end{align*}

	\section{Mechanische Wellen}
		\subsection{Energiedichte mechanischer Wellen}
			\[
				w = \frac{\mathrm{d}E}{\mathrm{d}V} = \frac12 \varrho \omega^2 \hat{y}^2
			\]
		
		\subsection{Massenbelag}
			\[
				m'=\frac{\Delta m}{\Delta l}
			\]	

		\subsection{Spannkraft}
			\[
				F=\frac{mc^2}{l}
			\]

			\begin{table}[h]
			\begin{tabular}{lll}
			$F\dots$ Spannkraft & $m\dots$ Saitenmasse & $l\dots$ Saitenlänge\\
			\end{tabular}
			\end{table}

	\section{elektrische Stromdichte $J$}
		\[ J = n e v_{\mathrm{Drift}} \]
		
		\begin{table}[h]
		\begin{tabular}{lll}
		$n\dots$Dichte der Ladungsträger & $e\dots$Elementarladung & $v_{\mathrm{Drift}}\dots$ Driftgeschwindigkeit\\
		\end{tabular}
		\end{table}
	
	\section{Schallwellen}
		\subsection{Schallpegel}
			\subsubsection{Pegeldifferenz}
				\[
					\Delta L=L_1-L_2=10\log\left(\frac{I_1}{I_2}\right)\mathrm{dB}
				\]

				\begin{table}[h]
				\begin{tabular}{lll}
				$L_{1,2}\dots$ Schallpegel & $I_{1,2}\dots$ Schallintensitäten & $L_1>L_2$\\
				\end{tabular}
				\end{table}

\clearpage

		\subsection{Dopplereffekt}

			\begin{table}[here]
			\begin{tabular}{ccl}
				Quelle & Beobachter & beobachtete Frequenz\\
				\toprule
				$\bullet$ & $\leftarrow \bullet$ & $f_{\mathrm{B}}=f_{\mathrm{Q}}\left(1+\frac{v_{\mathrm{B}}}{c}\right)$\\
				\midrule
				$\bullet$ & $\bullet \rightarrow$ & $f_{\mathrm{B}}=f_{\mathrm{Q}}\left(1-\frac{v_{\mathrm{B}}}{c}\right)$\\
				\midrule
				$\bullet \rightarrow$ & $\bullet$ & $f_{\mathrm{B}}=\frac{f_{\mathrm{Q}}}{1-\frac{v_{\mathrm{Q}}}{c}}$\\
				\midrule
				$\leftarrow \bullet$ & $\bullet$ & $f_{\mathrm{B}}=\frac{f_{\mathrm{Q}}}{1+\frac{v_{\mathrm{Q}}}{c}}$\\
				\midrule
				$\bullet \rightarrow$ & $\leftarrow \bullet$ & $f_{\mathrm{B}}=f_{\mathrm{Q}}\frac{c+v_{\mathrm{B}}}{c-v_{\mathrm{Q}}}$\\
				\midrule
				$\leftarrow \bullet$ & $\bullet \rightarrow$ & $f_{\mathrm{B}}=f_{\mathrm{Q}}\frac{c-v_{\mathrm{B}}}{c+v_{\mathrm{Q}}}$\\
				\midrule
				$\leftarrow \bullet$ & $\leftarrow \bullet$ & $f_{\mathrm{B}}=f_{\mathrm{Q}}\frac{c+v_{\mathrm{B}}}{c+v_{\mathrm{Q}}}$\\
				\midrule
				$\bullet \rightarrow$ & $\bullet \rightarrow$ & $f_{\mathrm{B}}=f_{\mathrm{Q}}\frac{c-v_{\mathrm{B}}}{c-v_{\mathrm{Q}}}$\\
				\bottomrule
			\end{tabular}
			\end{table}


			\begin{table}[h]
			\begin{tabular}{ll}
			$f_{\mathrm{B}}\dots$ Frequenz beim Beobachter & $f_{\mathrm{Q}}\dots$ Frequenz der Quelle\\
			$v_{\mathrm{B}}\dots$ Beobachtergeschwindigkeit relativ zur Luft & $v_{\mathrm{Q}}\dots$ Quellengeschwindigkeit relativ zur Luft\\
			$c\dots$ Schallgeschwindigkeit\\
			\end{tabular}
			\end{table}

		\subsection{Frequenzverschiebung elektromagnetischer Wellen}
			\begin{align*}
				f_{\mathrm{B}} &= f_{\mathrm{Q}}\sqrt{\frac{c+v}{c-v}}= f_{\mathrm{Q}}\sqrt{\frac{1+\beta}{1-\beta}} & \text{Bei Annäherung}\\
				f_{\mathrm{B}} &= f_{\mathrm{Q}}\sqrt{\frac{c-v}{c+v}}= f_{\mathrm{Q}}\sqrt{\frac{1-\beta}{1+\beta}} & \text{Bei Entfernung}
			\end{align*}

		\subsection{Schallintensitätspegel $L$}
			\begin{align*}
				L &= 10 \log{\frac{S}{S_0}}\mathrm{dB}  & S_1 &= S_0 10^{\frac{L}{10\mathrm{dB}}} & \frac{S_2}{S1} &= \left(\frac{r_1}{r_2} \right)^2
			\end{align*}
			Hinweis: diese Gleichungen gelten nur für Wellen mit ebenen Wellenfronten\newline

			\begin{table}[h]
			\begin{tabular}{ll}
			$S\dots$ Schallintensität & $S_0\dots$ Normschallintensität $=10^{-12}\frac{\mathrm{W}}{\mathrm{m}^2}$\\
			Leistung bei Kugeloberfläche $ P = 4\pi r^2I$
			\end{tabular}
			\end{table}

		\subsection{Machscher Kegel}
			\[ \sin \alpha = \frac{c_{\mathrm{Schall}}}{v_{\mathrm{Quelle}}}\] 
			
		\subsection{Schallwellengeschwindigkeit}
			\[ v = \frac{\hat{p}}{\varrho c}\cos(\omega t - kx) \]
		
		\subsection{Schallwellenwiderstand}
			oder auch Schallkennimpedanz
			\begin{align*}
				Z = \varrho c = \frac{\hat{p}}{\hat{v}} = \frac{\hat{p}^2}{I}
			\end{align*}
			
			\begin{table}[h]
			\begin{tabular}{lll}
			$Z\dots$ Schallkennimpedanz & $\hat{v}\dots$ Schnellenamplitude & $\hat{p}\dots$ Schallwechseldruck\\
			$\varrho\dots$ Dichte & $I\dots$ Intensität der Quelle\\
			\end{tabular}
			\end{table}

\clearpage

		\subsection{Reflexion, Transmission, Absorption}
			\begin{align*}
				\varrho_{\mathrm{s}} &= \left(\frac{Z_2-Z_1}{Z_2+Z_1}\right)^2 
				& \tau_{\mathrm{s}} &= \frac{I_{\mathrm{t}}}{I_{\mathrm{e}}}= \frac{4Z_1Z_2}{(Z_1+Z_2)^2} 
				& \alpha_{\mathrm{s}} &= \frac{I_{\mathrm{a}}}{I_{\mathrm{e}}}
			\end{align*}

			\begin{table}[h]
			\begin{tabular}{lll}
			$\varrho_{\mathrm{s}}\dots$ Reflexionsgrad & $\tau_{\mathrm{s}}\dots$ Transmissionsgrad & $\alpha_{\mathrm{s}}\dots$ Absorptionsgrad\\
			\end{tabular}
			\end{table}
		
	\section{Schnellenamplitude}
		\[ \hat{v} = \frac{\hat{p}}{\varrho c } = \frac{\hat{p}}{Z} \]

	\section{Auslenkung}
		\begin{align*}
			y &= \hat{y} \sin(\omega t - kx) & \text{mit} && \hat{y} &= \frac{\hat{p}}{\omega \varrho c} = \frac{\hat{p}}{Z \omega}
		\end{align*}	

	\section{Energiedichte $\omega$ bei el. mag. Wellen}
		\begin{align*}
			\omega &= \frac12 ( \epsilon_r \epsilon_0 E^2 + \mu_r \mu_o H^2) &\Rightarrow && S &= \omega c = \frac12 \epsilon_r \epsilon_0 \hat{E}^2 c =  \frac12 \mu_r \mu_0 \hat{H}^2 c
		\end{align*}
		$c\dots$Lichtgeschwindigkeit

	\section{Wellenwiderstand $Z$ (el/mag)}
		\[Z = \frac{E}{H} \]

		\begin{table}[h]
		\begin{tabular}{ll}
		$E\dots$ Elektrisches Feld & $H\dots$ Magnetfeld\\
		\end{tabular}
		\end{table}

	\section{2 getrennte Leiter}
		\begin{align*}
			C' &= \frac{\epsilon_r \epsilon_0 \pi}{\ln(\frac{a}{r}}) & L' &= \frac{\mu_r \mu_0}{\pi} \cdot \ln(\frac{a}{r})
		\end{align*}

		\begin{table}[h]
		\begin{tabular}{ll}
		$a\dots$Abstand der Leiter & $r\dots$Durchmesser eines Leiters\\
		\end{tabular}
		\end{table}
		
	\section{2 Leiter ineinander}
		\begin{align*}
			C' &= \frac{2\epsilon_r \epsilon_0 \pi}{\ln(\frac{D}{d}}) & L' &= \frac{\mu_r \mu_0}{2\pi} \cdot \ln(\frac{D}{d})
		\end{align*}

		\begin{table}[h]
		\begin{tabular}{ll}
		$d\dots$Durchmesser des Innenleiters & $D\dots$Durchmesser des Außenleiters\\
		\end{tabular}
		\end{table}
	\section{Ausbreitungsgeschwindigkeit bei Wellen}

		\begin{table}[here]
			\begin{tabular}{l l}
				Longitudinalwellen in Gasen & $ c = \sqrt{\frac{\kappa p}{\varrho}} $\\\midrule
				Longitudinalwellen in Flüssigkeiten & $ c = \sqrt{\frac{K}{\varrho}} $\\\midrule
				Longitudinalwellen in Stäben & $ c = \sqrt{\frac{E}{\varrho}} $\\\midrule
				Torionswellen in Rundstäben & $ c = \sqrt{\frac{G}{\varrho}} $ \\\midrule
				El./Mag. Wellen im Vakuum & $c = \frac{1}{\sqrt{\epsilon_0 \mu_0}} $\\\midrule
				El./Mag. Wellen in Materie	 	& $ c = \frac{1}{\sqrt{\epsilon_0 \epsilon_r \mu_0 \mu_r}} $\\\midrule
				Auf Saiten					& $ c = \sqrt{\frac{F}{A \varrho}} $
			\end{tabular}
			\caption{Ausbreitungsgeschwindigkeit}
		\end{table}
		
		\begin{table}[h]
		\begin{tabular}{lll}
		$\kappa\dots$ Isotropenexponent & $\varrho\dots$ Dichte & $p\dots$ Druck\\
		$A\dots$ Querschnitt des Drahts & $F\dots$ Spannkraft & $K\dots$ Kompressionsmodul\\
		$E\dots$ Elastizitätsmodul & $G\dots$ Schubmodul\\
		\end{tabular}
		\end{table}

\clearpage

	\section{Wellenwid. im Freiraum und Leitungen}

		\begin{table}[here]
			\begin{tabular}{l c c}
				& Leitungen & Freiraum\\\toprule
				Definition & $Z_L = \frac{\underline{U}}{\underline{I}}$ & $Z_F = \frac{\underline{E}}{\underline{H}} $\\\midrule
				Verlustbehaftet & $Z_L = \sqrt{\frac{R' + j\omega L'}{G'+j\omega C'}}$ & $ Z_L  = \sqrt{\frac{\mu_r \mu_0}{\epsilon_r \epsilon_0 - j\frac{\kappa}{\omega}}} $ \\\midrule
				Verlustlos & $Z_L = \sqrt{\frac{L'}{C'}}$ &$Z_L = \sqrt{\frac{\mu_0}{\epsilon_o}} = 376,6 \Omega$\\\bottomrule
			\end{tabular}
			\caption{Wellenwiderstand}
		\end{table}
	
		\begin{table}[h]
		\begin{tabular}{lll}	
		$R' = \frac{R}{l} \dots$ längs & $C' = \frac{C}{l} \dots$ & $L' = \frac{L}{l} \dots$ \\
		$G' = \frac{G}{l} \dots$ quer & $l \dots$ länge & $\kappa \dots$ Spezifische Leitfähigkeit\\
		\end{tabular}
		\end{table}



\chapter{Interferenz}
	\section{Interferenz bei mechanischen Wellen}
		\subsection{konstruktive Interferenz}
			Gangunterschied $\Delta$
			\[\Delta = n \lambda \]
			Phasenverschiebung $\Delta\varphi$
			\[\Delta\varphi = n 2 \pi \]
			$n = 0,1,2,3,\dots$
			
		\subsection{destruktive Interferenz}
			Gangunterschied $\Delta$
			\[\Delta = (2n + 1) \frac{\lambda}{2} \]
			Phasenverschiebung $\Delta\varphi$
			\[\Delta\varphi = (2 n +1)\pi \]
			$n = 0,1,2,3,\dots$		
		
	\subsection{Allgemeiner Bezug Phasenverschiebung / Gangunterschied}
		\[ \frac{\varphi}{2\pi} = \frac{\Delta}{\lambda} \]
		
	\section{Interferenz bei Schallwellen}
		(gilt bei gleicher Amplitude, Frequenz, Ausbreitungsrichtung
			
		\subsection{konstruktive Interferenz}
			Gangunterschied $\Delta$:
			\[ \Delta = n \lambda \]
			$n = 0,1,2,3,\dots$
			
		\subsection{destruktive Interferenz}
			Gangunterschied $\Delta$
			\[ \Delta = (2n +1) \frac{\lambda}{2} \]
			$n = 0,1,2,3,\dots$
			
			Außerdem gilt für weit entfernte Punkte:
			\[\sin \alpha = \frac{\Delta}{g} \]
			
			\begin{table}[h]
			\begin{tabular}{ll}
			$\Delta\dots$Gangunterschied & $g\dots$Länge zwischen den beiden Quellen\\
			\end{tabular}
			\end{table}
		
	\section{Interferenz bei Lichtwellen}
		\subsection{Am engen Spalt}
			\begin{align*}
				\sin\alpha_{\min} &= n \frac{\lambda}{b} & \sin\alpha_{\max} = \left(n + \frac12\right) \frac{\lambda}{b}
			\end{align*}
			
			\begin{table}[h]
			\begin{tabular}{ll}
			$n=1,2,3,\dots$ & $b\dots$ Spaltbreite\\
			\end{tabular}
			\end{table}
			
		\subsection{Am Doppelspalt}
			\begin{align*}
				\sin\alpha_{min} &= \left(n + \frac12\right) \frac{\lambda}{d} & \sin\alpha_{max} = n\frac{\lambda}{d}
			\end{align*}
			\begin{table}[h]
			\begin{tabular}{ll}
			$n=0,1,2,3,\dots$ & $d\dots$Breite von Spaltmitte 1 bis Spaltmitte 2\\
			\end{tabular}
			\end{table}
			
		\subsection{Beugungsgitter}
			\begin{align*}
				\sin\alpha_{min} &= \left(n + \frac12\right) \frac{\lambda}{g} 
				& \sin\alpha_{max}&= n\frac{\lambda}{g} 
				& s&=\frac{1}{g}
			\end{align*}
			
			\begin{table}[h]
			\begin{tabular}{lll}
			$n=0,1,2,3,\dots$ & $g\dots$ Gitterkonstante & $s\dots$ Linienanzahl pro Fläche\\
			\end{tabular}
			\end{table}

			\subsubsection{Auflösung}
				\begin{align*}
					\frac{\lambda}{\Delta\lambda}&=mp & \frac{\lambda}{\Delta\lambda}&=msb
					& \Delta x&=pg=\frac{\lambda}{\Delta \lambda}\frac{1}{ms}
				\end{align*}
				
				\begin{table}[h]
				\begin{tabular}{lll}
				$m\dots$ Ordnung & $p\dots$ Linienzahl & $b\dots$ Gitterbreite\\
				\end{tabular}
				\end{table}
		
			\subsubsection{Abstand zweier Maxima-Linien}
				\[
					\Delta s=\frac{l}{g}(\lambda_1-\lambda_2)
				\]

				$l\dots$ Abstand Gitter zu Schirm

\chapter{Quantenoptik}
	\section{Photonenenergie}
		Im Grenzfall gilt $W_{\mathrm{Austritt}} = E_{\mathrm{Grenz}}$
		\[ E_{ph} = h f = \frac{hc}{\lambda} = \frac{h'}{\lambda} \]
		\begin{table}[h]
		\begin{tabular}{ll}
		$h=6,626\cdot10^{-34}$Js & $h'\dots$ 1,24$\mu$meV\\
		\end{tabular}
		\end{table}
		
	\section{Impuls $p$ eines Photons}
		\[p = \frac{hf}{c} = \frac{h}{\lambda} = m v\]
		
		Photonenmasse:
		\begin{align*}
			m &= \frac{E_{ph}}{c^2} = \frac{hf}{c^2} &&& E_{kin} &= E - E_0 = (m - m_0)c^2 \\
			v &= c \sqrt{1-\frac{1}{1+(\frac{eU_A}{m_o c^2})^2}} &&&
		\end{align*}
		
		$c\dots$ Lichtgeschwindigkeit, üblicherweise $3\cdot10^8\nicefrac{\mathrm{m}}{\mathrm{s}}$
		
	\section{Strahlungsleistung $\Phi$ eines Lichtstrahles}
		\begin{align*}
			\Phi &= E_{\mathrm{ph}} \dot{N}_{\mathrm{ph}} = IA 
			& \dot{N} &= \frac{\lambda}{hc} \Phi = \frac{\lambda}{h'} \Phi 
			& I_{\mathrm{ph}} &= e\dot{N} \\
			\Phi & = \Phi_0 \mathrm{e}^{-\alpha d} 
		\end{align*}
		
		\begin{table}[h]
		\begin{tabular}{lll}
		$I_{ph}\dots$ Strom & $\dot{N}\dots$ Photonenstrom & $I\dots$ Intensität $\left[\frac{W}{m^2}\right]$ \\
		$A\dots$ Fläche & $\alpha\dots$ Absorbtionskoeffizient & $d\dots$ Dicke des Kristalls\\
		$h'\dots$ 1,24$\mu$meV\\
		\end{tabular}
		\end{table}
		
	\section{Beugung von Lichtwellen}
		oder Heisenbergsche Unschärferelation genannt
		\begin{align*}
			e U_A &= \frac12 m v^2 & v &= \sqrt{\frac{2eU_A}{m}}
		\end{align*}
		Beugung bei kleinen Winkeln $\tan \alpha = \frac{\Delta px}{py} = \frac{\Delta px}{\frac{h}{\lambda}} $\newline
		Beugungsminimum: $\sin \alpha = \frac{\lambda}{b} = \frac{\lambda}{\Delta x}$ \newline
		Bei kleinen Winkeln: $\tan x  = \sin x $ \newline
		\begin{align*}
			\frac{\Delta px}{\frac{h}{\lambda}} &= \frac{\lambda}{\Delta x} & \Delta px \Delta x &= h
		\end{align*}

		\subsection{Brechung, Reflexion}
			\begin{align*}
				\varrho_{\mathrm{s}}&=\left(\frac{n_2-n_1}{n_2+n_1}\right)^2 
				& \tau_{\mathrm{s}}&=\frac{4n_1n_2}{(n_1+n_2)^2}=\frac{I_{\mathrm{t}}}{I_{\mathrm{e}}}
				& \alpha_{\mathrm{s}}&=\frac{I_{\mathrm{a}}}{I_{\mathrm{e}}}
			\end{align*}

			\begin{table}[h]
			\begin{tabular}{lll}
			$\varrho_{\mathrm{s}}\dots$ Reflexionsgrad & $\tau_{\mathrm{s}}\dots$ Transmissionsgrad & $\alpha_{\mathrm{s}}\dots$ Absorptionsgrad\\
			\end{tabular}
			\end{table}
			\begin{table}[h]
			\begin{tabular}{ll}
			Stoff & Brechzahl\\
			\toprule
			Wasser & $n_{\mathrm{H_2O}}=1,33$\\
			Glas & $n_{\mathrm{Glas}}=1,5$\\
			Luft & $n_{\mathrm{Luft}}=1$\\
			\end{tabular}
			\caption{Häufige Brechzahlen}
			\end{table}

	\section{Elektronen im elektrischen Feld}
		\begin{align*}
			F_{\text{el}}&= -eE & a=-\frac{e}{m}E
		\end{align*}

		Arbeit: $W_{12}=eU_{12}=\frac{1}{2}m(v_2^2-v_1^2)$
		Wellenlänge des Elektrons:

		\begin{align*}
			E_{\text{kin}}&= eU_A & \frac{1}{2}mv^2&= eU_A & v=\sqrt{\frac{2eU_A}{m}}
		\end{align*}

		$h=6,626\cdot10{-34}$Js \\
		mit p=mv=$\nicefrac{h}{\lambda}$ folgt die Wellenlänge $\lambda=f(U_A)$
		\[
			\lambda=\frac{h}{\sqrt{2eU_Am}}
		\]
		\begin{table}[h]
		\begin{tabular}{ll}
		$U_A\dots$ Beschleunigungsspannung & $m\dots$ Elektronenmasse\\
		\end{tabular}
		\end{table}

	\section{Bahnkurve des Elektrons im E-Feld}
		\[
			a_x=0; a_y=\frac{e}{m}E; v_x=v_x;v_y=\frac{e}{m}Et
		\]
		\begin{align*}
			V_x&= v_xt & v_y&= \frac{1}{2}\frac{l}{m}Et^2 & y&= \frac{1}{2}\frac{e}{m}E\frac{x^2}{v^2} & y&= \frac{1}{4}\frac{U_K}{U_A}\frac{x^2}{d}
		\end{align*}
		\begin{table}[h]
		\begin{tabular}{ll}
		$d\dots$ Plattenabstand & $e=1,602\cdot10^{-19}$As\\
		\end{tabular}
		\end{table}
	
	\section{Bewegung von Elektronen im Magnetfeld}
		\[
			F_{\text{L}}=qvB \Rightarrow e^-:F_{\text{L}}=-evB
		\]
		\[
			F_{\mathrm{L}}=F_{\mathrm{Zp}} \Rightarrow qvB=\frac{mv^2}{r}
		\]
		\subsection{Linke-Hand-Regel}
			\begin{tabularx}{\textwidth}{l l}
				Daumen & Elektronenflussrichtung\\
				Zeigefinger & Magnetfeldrichtung\\
				Mittelfinger & Lorentzkraft\\
			\end{tabularx}
        
		\subsection{Rechte-Hand-Regel}
			\begin{tabularx}{\textwidth}{l l}
				Daumen & technische Stromrichtung\\
				Zeigefinger & Magnetfeldrichtung\\
				Mittelfinger & Lorentzkraft\\
			\end{tabularx}
	
	\section{Radius im Fadenstrahlrohr}
		\begin{align*}
			\frac{mv^2}{r}&= evB \Rightarrow r = \frac{mv}{eB} & T= \frac{2\pi r}{v} &= \frac{2\pi m}{eB}
		\end{align*}

		\begin{align*}
			T &= \frac{2\pi}{B}\frac{m}{e} & \omega &= \frac{2\pi}{T}=\frac{eB}{m}
		\end{align*}

		\subsection{Zyklotron}
			\begin{align*}
				\tau &= \frac{\mu m^*}{e} 
				& v_\text{{th}} &= \sqrt{\frac{3kT}{m^*}} 
				& l &= v_{th}\tau 
				& f_{\text{res}} &= \frac{1}{\tau} 
				& B_{\text{res}} &= \frac{2\pi f_{\text{res}} m^*}{e}
			\end{align*}

			\begin{table}[h]
			\begin{tabular}{ll}
			$\tau\dots$ Mittlere Flugdauer(Streuzeit), Relaxationszeit & $m^*\dots$ Reduzierte Masse\\
			$l\dots$ mittlere freie Weglänge & $v_{\text{th}}\dots$ thermische Geschwindigkeit \\
			$f_{\text{res}}\dots$ Zyklotronenresonanzfrequenz & $B\dots$ Magnetfeld, [T]\\
			\end{tabular}
			\end{table}

	\section{Leiterschleife im Magnetfeld}
		\begin{align*}
			M&=IBA\sin\alpha=IBab\sin\alpha & ab&=A
		\end{align*}

	\section{Relativistische Bewegung}
		\begin{align*}
			E &= mc^2\text{(ruhend)} & \text{bzw. } E_0 &= m_0c^2 & m &= \frac{m_0}{\sqrt{1-\frac{v^2}{c^2}}}
		\end{align*}

	\section{Elektronen im Atom}
		Kreisfrequenz eines Elektrons $\Rightarrow F_{\text{Z}}=F_{\text{coulomb}}$ 
		wobei $Q=q_1q_2=Ze^2$, allerdings nur bei Atomen, die nur ein Außenelektron haben.
		\begin{align*}
			\omega^2 &= \frac{Ze^2}{4\pi\varepsilon_0m_0r^3} 
			& E_{\text{kin}}=\frac{1}{2}J\omega^2=\frac{1}{2}mr^2\omega^2=\frac{Ze^2}{8\pi\varepsilon_0}\frac{l}{r}
		\end{align*}
		\begin{align*}
			E_{\text{pot}} &= -\frac{1}{4\pi\varepsilon_0}\frac{Ze^2}{r} 
			& E_{\text{ges}} &= E_{\text{kin}}+E_{\text{pot}} = -\frac{Ze^2}{8\pi\varepsilon_0}\frac{l}{r}
		\end{align*}
	
		\begin{table}[h]
		\begin{tabular}{ll}
		$Z\dots$ Ordnungszahl\\
		\end{tabular}
		\end{table}
	\section{Bohrsches Postulat}
		für stehende Wellen gilt:
		\[
			\text{Bahnumfang} = 2\pi r = n\lambda = n\lambda = n\frac{h}{p} \Rightarrow l = rp = n\frac{h}{2\pi}
		\]
		\[
			l = n\frac{h}{2\pi} = n\bar{h} = m_0r\omega
		\]

		mit $\omega^2=\dots$ ergibt sich:
		\[
			r_n=\frac{4\pi\varepsilon_0\bar{h}^2}{e^2m_0}\frac{n^2}{Z}
		\]
		mit $e^2m_0=5,29\cdot10^{-11}$m.

		Die Energie eines Elektrons auf einer Bahn
		\[
			E_n=-Z^2\frac{e^4m_0}{32\pi^2\varepsilon_0^2\bar{h}^2}\frac{1}{n^2}
		\]
		mit $32\pi^2\varepsilon_0^2\bar{h}^2=13,6$eV.

		\subsection{Bindungs- oder Ionisierungsenergie}
		$\frac{e^4m_0}{32\pi^2\varepsilon_0^2\bar{h}^2}$ ist ein konstanter Faktor = 13,6 eV, gegebenenfalls 
		mit reduzierter Masse multiplizieren oder durch $\varepsilon^2$ dividieren.

	\section{Fermi-Energie}
		\[
			E_{\text{F}}=\frac{\bar{h}^2}{2m}(3\pi^2)^{\frac{2}{3}}n^{\frac{2}{3}} \text{ mit } n=\frac{N}{V}
		\]
        
		Übergang zur Fermi-Geschwindigkeit:
		\[
			E_{\text{F}}=\frac{\bar{h}k_f}{2m}=\frac{p_{\text{F}}^2}{2m}=\frac{1}{2}mv_{\text{F}}^2
		\]

		\begin{table}[h]
		\begin{tabular}{lll}
		$n\dots$ Ladungsträgerdichte & $N\dots$ Anzahl der Ladungsträger & $\bar{h}=1,055\cdot10^{-34}\mathrm{Js}=6,58\cdot10^{-16}\mathrm{eVs}$\\
		\end{tabular}
		\end{table}
        
		\[
			v_{\text{F}}=\sqrt{\frac{2E_{\text{F}}}{m}}
		\]
       
       		\begin{table}[h]
		\begin{tabular}{ll} 
		$v_{\text{F}}\dots$ Fermigeschwindigkeit & $k_{\text{F}}\dots$ Wellenzahl, $k=\frac{2\pi}{\lambda}$\\
		\end{tabular}
		\end{table}
        
		\[
			v_{\text{F}}=\frac{\bar{h}}{m}(3\pi^2n)^{\frac{1}{3}}
		\]

	\section{Zustandsdichte bei $T=0$K}
		\[
			D(E)=\frac{1}{2\pi^2}\left(\frac{2m^{\frac{2}{3}}}{\bar{h}}\right)E^{\frac{1}{2}}
		\]

		Mit zunehmender Temperatur gilt die Besetzungswahrscheinlichkeit
		\[
			f(E)=\frac{1}{1+e^{\frac{E-E_{\text{F}}}{kT}}}
		\]
	
	\section{Drehmoment im homogenen elektrischen Feld}
		\begin{align*}
			M &= Fr\sin\beta & \vec{M}&=\vec{p}\times\vec{E} & \vec{p}&= q\vec{r} \text{ el. Dipolmoment}
		\end{align*}

	\section{Drehmoment im magnetischen Feld}
		\begin{align*}
			\vec{m}&=I\vec{A} \text{ mag. Dipolmoment} & m&=N\mu_B \text{ mag. Dipolmoment}
		\end{align*}

		\begin{table}[h]
		\begin{tabular}{ll}
		$\mu_B=9,274\cdot10^{-24}\text{Am}^2\dots$ Bohrsches Magneton & $N\dots$ Ladungsträger\\
		\end{tabular}
		\end{table}

\chapter{Halbleiter}
	\section{Eigenleitung}
		\begin{align*}
			n_i &= n_{i0}T^{\frac{3}{2}}e^{-\frac{E_\text{g}}{2kT}} & \kappa &=en_i(\mu_n+\mu_p) & n_i &= p=n
		\end{align*}
	
	\section{Temperaturabhängigkeit der Beweglichkeit $\mu$}
		\begin{align*}
			\mu(T) &= \mu_0\left(\frac{T}{T_0}\right)^{\frac{3}{2}} & R(T) &= R_0e^{-\frac{E_{\text{g}}}{2kT}}
		\end{align*}
	
	\section{Temperaturabhängigkeit des Widerstands}
		\begin{align*}
			\frac{R_1}{R_2}&=e^{\frac{E_{\mathrm{A}}}{2k}\left(\frac{1}{T_1}-\frac{1}{T_2}\right)} 
			& E_{\mathrm{A}}&=\frac{2k}{\frac{1}{T_1}-\frac{1}{T_2}}\cdot\ln\frac{R_1}{R_2}
		\end{align*}
		\begin{table}[h]
		\begin{tabular}{ll}
		n-Typ & p-Typ\\
		\toprule
		$R\sim e^{\frac{E_{\mathrm{D}}}{2kT}}$ & $R\sim e^{\frac{E_{\mathrm{A}}}{2kT}}$\\
		\end{tabular}
		\end{table}

	\section{Proportionalitäten}
		\begin{align*}
			n\sim e^{-\frac{E_{\mathrm{D}}}{2kT}} && p\sim e^{-\frac{E_{\mathrm{A}}}{2kT}} && R\sim e^{E_{\mathrm{D}}{2kT}}
		\end{align*}

	\section{Ladungsträgerdichte $n$, $p$}
		\[
			n=\frac{m}{V}\frac{N_\text{A}}{M}
		\]

		\begin{table}[h]
		\begin{tabular}{lll}
		$\nicefrac{m}{V}=\rho\dots$ Dichte & $\nicefrac{N_{\text{A}}}{M}=m_{\text{Atom}}\dots$ Atommasse & $m\dots$ Masse\\
		$V\dots$ Volumen & $N_{\text{A}}\dots$ Avogadrokonstante, $6,022\cdot10^{23}\nicefrac{1}{\mathrm{mol}}$ & $M\dots$ Molare Masse in $\nicefrac{g}{\mathrm{mol}}$\\
		\end{tabular}
		\end{table}

\clearpage

	\section{Intrinsische Trägerdichte $n_i$}
		\subsection{n-Typ}
		aus der V. bis VII. Hauptgruppe
			\begin{align*}
				n &= \frac{n_{\text{D}}}{2}+
				\sqrt{\left(\frac{n_{\text{D}}}{2}\right)^2+n_i^2} 
				& np &=n_i^2 & p &=\frac{n_i^2}{n}
			\end{align*}

			\begin{table}[h]
			\begin{tabular}{ll}
			Für $n_{\text{D}}\gg n_i \rightarrow n=n_{\text{D}}$ & Für $n_{\text{D}}\ll n_i \rightarrow n=n_i$\\
			\end{tabular}
			\end{table}

		\subsection{p-Typ}
		aus der I. bis III. Hauptgruppe
			\begin{align*}
				p &= \frac{n_{\text{A}}}{2}+
				\sqrt{\left(\frac{n_{\text{A}}}{2}\right)^2+n_i^2} 
				& np &= n_i^2 & n &= \frac{n_i^2}{p}
			\end{align*}

			\begin{table}[h]
			\begin{tabular}{ll}
			Für $n_{\text{A}}\gg n_i \rightarrow n=n_{\text{A}}$ & Für $n_{\text{A}}\ll n_i \rightarrow n=n_i$\\
			\end{tabular}
			\end{table}

	\section{Spezifische Leitfähigkeit $\kappa$}
		\[
			\kappa=e(n\mu_{\text{n}}+p\mu_{\text{p}}), [\kappa]=\nicefrac{1}{\Omega\text{cm}}
		\]

		\begin{table}[h]
		\begin{tabular}{ll}
		Für starke n-Dotierung $\kappa_{\text{n}}=en_{\text{D}}\mu_{\text{n}}$ & Für starke p-Dotierung $\kappa_{\text{p}}=en_{\text{A}}\mu_{\text{p}}$\\
		\end{tabular}
		\end{table}
	
	\section{Störstellenreserve}
		\begin{tabularx}{\textwidth}{c c}
		für n-Leiter & für p-Leiter\\
		$n=\sqrt{\frac{n_{\text{D}}N_{\text{L}}}{2}e^{\frac{E_{\text{D}}}{2kT}}}$
		& $n=\sqrt{\frac{n_{\text{A}}N_{\text{L}}}{2}e^{\frac{E_{\text{A}}}{2kT}}}$ \\
		\end{tabularx}

		\begin{table}[h]
		\begin{tabular}{ll}
		$N_{\text{L}}\dots$ Zustandsdichte Leitungsband & $E_{\text{D}}\dots$ Donatoren-Energielücke\\
		$E_{\text{A}}\dots$ Akzeptoren-Energielücke\\
		\end{tabular}
		\end{table}
	
	\section{Räumliche Ausbreitung einer Störung}
		\begin{align*}
			\frac{\mathrm{d}N}{\mathrm{d}t}=\dot{N}=-\left(\frac{\mathrm{d}\Delta n}{\mathrm{d}t}\right)\cdot A\cdot D
		\end{align*}

		\begin{table}[h]
		\begin{tabular}{lll}
		$\dot{N}\dots$ Ladungsträgerstrom & $A\dots$ Fläche & $D\dots$ Diffusionskoeffizient\\
		\end{tabular}
		\end{table}
	
	\section{Diffusionsspannung}
		\begin{align*}
			U_{\mathrm{d}}&=\frac{kT}{e}\cdot\ln\left(\frac{n_{\mathrm{D}}\cdot n_{\mathrm{A}}}{n_{\mathrm{i}}^2}\right)
			& U_{\mathrm{dR}}&=25,9\mathrm{mV}\cdot\ln\left(\frac{n_{\mathrm{D}}\cdot n_{\mathrm{A}}}{n_{\mathrm{i}}^2}\right)
		\end{align*}
		
		\begin{table}[h]
		\begin{tabular}{ll}
		$kT\dots$ 25,9meV bei Raumtemperatur (300K) & $U_{\mathrm{dR}}\dots$ Diffusionsspannung bei Raumtemperatur\\
		\end{tabular}
		\end{table}

		\subsection{Entfernung vom pn-Übergang}
			\subsubsection{Allgemein}
				\[
					\frac{\partial\Delta n}{\partial t}=\frac{D_{\mathrm{n}}\cdot\partial^2\Delta n}{\partial x^2}-\frac{\Delta n}{\tau_{\mathrm{n}}}
				\]

			\subsubsection{Stationärer Zustand}
				\begin{align*}
					L_{\mathrm{n}}&=\sqrt{D_{\mathrm{n}}\tau}\cdot5\approx x_{\mathrm{n}}
					& L_{\mathrm{p}}&=\sqrt{D_{\mathrm{p}}\tau}\cdot5\approx x_{\mathrm{p}}
				\end{align*}

				$x_{\mathrm{{n,p}}}\dots$ ungefähre Entfernung vom pn-Übergang\\

	\section{Driftgeschwindigkeit}
		\begin{align*}
			v_{\text{d}0} &= -\frac{e}{m}\tau E_0 & v_{\text{d}0} &=-\mu E_0 & \tau &= \frac{\kappa m_0}{e^2n} 
			& j &= -env_{\text{d}0}=\kappa E_0
		\end{align*}

		\begin{table}[h]
		\begin{tabular}{ll}
		$E_0\dots$ Elektrische Feldstärke & $e=1,602\cdot10^{-19}$As\\
		\end{tabular}
		\end{table}

	\section{Ionisierungswahrscheinlichkeit}
		\[
			f_{\text{I}}=1-f(E_{\text{D}})
		\]

		$E_0\dots$ Besetzungswahrscheinlichkeit des Donatoratoms

	\section{Effektive Zustandsdichte der Energiebänder}
		\begin{align*}
			n &= N_{\text{L}}e^{-\frac{E_{\text{L}}-E_{\text{F}}}{kT}} 
			& p &= N_{\text{v}}e^{-\frac{E_{\text{F}}-E_{\text{V}}}{kt}} 
			& n_{\text{i}}^2 = np = N_{\text{L}}N_{\text{V}}e^{-\frac{E_{\text{g}}}{kT}}
		\end{align*}

\clearpage

	\section{Austrittsarbeit}
		\begin{align*}
			W_{\mathrm{A}}&=E_{\mathrm{ph}}=\frac{hc}{\lambda}=\frac{h'}{\lambda} & E_{\mathrm{kin}}&=\frac{1}{2}mv^2=E_{\mathrm{ph}}-W_{\mathrm{A}}
		\end{align*}
		
		\begin{table}[h]
		\begin{tabular}{lll}
		$E_{\mathrm{ph}}\dots$ Photonenenergie & $W_{\mathrm{A}}\dots$ Austrittsarbeit & $h'\dots$ 1,24$\mu$meV \\
		\end{tabular}
		\end{table}
	
	\section{Shockley-Gleichung}
		\begin{align*}
			I&=I_{\mathrm{s}}\left(e^{\frac{eU}{kT}}-1\right) 
			& I_{\mathrm{R}}&=I_{\mathrm{s}}\left(e^{\frac{U}{25,9\mathrm{mV}}}-1\right) 
			& I_{\mathrm{s}}&=Ae\left(n_{\mathrm{p}}\frac{D_{\mathrm{n}}}{L_{\mathrm{n}}}+p_{\mathrm{n}}\frac{D_{\mathrm{p}}}{L_{\mathrm{p}}}\right)
		\end{align*}
	
		\begin{table}[h]
		\begin{tabular}{ll}
		$I_{\mathrm{s}}\dots$ Sperrsättigungsstrom & $L_{\mathrm{n,p}}\dots$ Diffusionslänge der Ladungsträger\\
		$A\dots$ Diodenfläche & $e=1,602\cdot10^{-19}$As\\
		$I_{\mathrm{R}}\dots$ Strom bei Raumtemperatur & \\
		\end{tabular}
		\end{table}
	
	\section{Differenzieller Widerstand $r$}
		\begin{align*}
			r&=\frac{\mathrm{d}U}{\mathrm{d}I} 
			& \frac{1}{r}&=\frac{\mathrm{d}I}{\mathrm{d}U}=I_{\mathrm{s}}\cdot\frac{e}{kT}\cdot e^{\frac{eU}{kT}}
			& \frac{1}{r_{\mathrm{R}}}&=\frac{\mathrm{d}I}{\mathrm{d}U}=I_{\mathrm{s}}\cdot\frac{1}{25,9{\mathrm{mV}}}\cdot e^{\frac{U}{25,9\mathrm{mV}}}
		\end{align*}

\clearpage

	\section{Eigenschaften wichtiger Halbleiter}
		\begin{table}[here]
		\begin{tabularx}{\textwidth}{llll}
		& Ge & Si & GaAs\\
		\toprule
		Kristallstruktur & Diamant & Diamant & Zinkblende\\
		\midrule
		Gitterkonstante $a$ in $10^{-10}$m & 5,65771 & 5,43043 & 5,65325\\
		\midrule
		lin. Ausdehnungskoeff. $\alpha$ in $10^{-6}\mathrm{K}^{-1}$ & 5,90 & 2,56 & 6,86\\
		\midrule
		spez. Wärmekapazität $c$ in $\nicefrac{\mathrm{kJ}}{\mathrm{kgK}}$ & 0,31 & 0,70 & 0,35\\
		\midrule
		Wärmeleitfähigkeit $\lambda$ in $\nicefrac{\mathrm{W}}{\mathrm{mK}}$ & 64 & 145 & 46\\
		\midrule
		Schmelzpunkt $\vartheta_{\mathrm{S}}$ in °C (K) & 937 (1210,15) & 1415 (1688,15) & 1238 (1511,15)\\
		\midrule
		Atomdichte $\nicefrac{\mathrm{N}}{\mathrm{V}}$ in $10^{22}\mathrm{cm}^{-3}$ & 4,42 & 5,0 & 4,42\\
		\midrule
		Dichte $\varrho$ in $\nicefrac{\mathrm{kg}}{\mathrm{m}^3}$ & 5326,7 & 2328 & 5320\\
		\midrule
		Molmasse $M$ in $\nicefrac{\mathrm{g}}{\mathrm{mol}}$ & 72,60 & 28,09 & 144,63\\
		\midrule
		Bandgap $E_{\mathrm{g}}$ in eV & 0,660 & 1,11 & 1,43\\
		\midrule
		intrinsische Trägerdichte $n_{\mathrm{i}}$ in $\mathrm{cm}^{-3}$ & $2,33\cdot10^{13}$ & $1,02\cdot10^{10}$ & $2,00\cdot10^6$\\
		\midrule
		Effektive Zustandsdichte & &\\
		im Leitungsband $N_{\mathrm{L}}$ in $\mathrm{cm}^{-3}$ & $1,24\cdot10^{19}$ & $2,85\cdot10^{19}$ & $4,55\cdot10^{17}$\\
		im Valenzband $N_{\mathrm{V}}$ in $\mathrm{cm}^{-3}$ & $5,35\cdot10^{18}$ & $1,62\cdot10^{19}$ & $9,32\cdot10^{18}$\\
		\midrule
		Beweglichkeit & &\\
		$\mu_{\mathrm{n}}$ in $\nicefrac{\mathrm{cm}^2}{\mathrm{Vs}}$ & 3900 & 1350 & 8500\\
		$\mu_{\mathrm{p}}$ in $\nicefrac{\mathrm{cm}^2}{\mathrm{Vs}}$ & 1900 & 480 & 435\\
		\midrule
		relative Dielektrizitäts-/Permittivitätszahl & 16 & 11,8 & 12,9\\
		\midrule
		Lebendsdauer d. $e^-$ $\tau$ (Größenordnung) & ms & $\mu$s & ns\\
		\bottomrule
		\end{tabularx}
		\label{wichtige Halbleiter}
		\caption{Eigenschaften wichtiger Halbleiter, Quelle: HMS, 12. Auflage}
		\end{table}

\chapter{Hall-Effekt}
	\begin{table}[h]
	\begin{tabular}{ll}
	$b\dots$ Breite & $d\dots$ Dicke \\
	$B\dots$ Flussdichte & $v_{\text{D}}\dots$ Driftgeschwindigkeit der Elektronen \\
	$n\dots$ Ladungsträgerdichte, ggf. $N_{\text{A}}$ oder $N_{\text{D}}$ & $N\dots$ Teilchenanzahl\\
	$V\dots$ Volumen\\
	\end{tabular}
	\end{table}

	\section{Kräftegleichgewicht}
		\[
			F_{\text{el}}=bE_{\text{H}}=e\frac{U_{\text{H}}}{b}=F_{\text{L}}
		\]
	
	\section{Hallspannung}
		\[
			U_{\text{H}}=v_{\text{D}}bB=\frac{IB}{end}=\frac{R_{\text{H}}IB}{d}
		\]
	
	\section{Hallkoeffizient $R_{\text{H}}\left[\frac{\text{m}^3}{\text{As}}\right]$}
		\[
			R_{\text{H}}=\frac{1}{en}=\frac{\mu}{\kappa}=\frac{U_{\text{H}}d}{IB}
		\]

	\section{Längswiderstand}
		\[
			R=\frac{l}{\kappa bd}=\frac{l}{\kappa A}
		\]

	\section{Spezifische Leitfähigkeit $\kappa$}
		\[
			\kappa=\frac{l}{Rbd}=e(n\mu_{\text{n}}+p\mu_{\text{p}}), [\kappa]=\nicefrac{1}{\Omega\text{cm}}
		\]
	
	\section{Ladungsträgerdichte}
		\[
			n=\frac{N}{V}=\frac{1}{eR_{\text{H}}}
		\]
	
	\section{Beweglichkeit $\mu\left[\frac{\text{cm}^2}{\text{Vs}}\right]$}
		\[
			\mu=\kappa R_{\text{H}}=\frac{\kappa}{eR_{\text{H}}}
		\]

	\section{Widerstand im Magnetfeld}
		\[
			R=R_0(1+\alpha B^2)
		\]

		$R_0\dots$ Widerstand ohne Magnetfeld

	\section{Temperaturabhängigkeit der Leitfähigkeit}
		\[
			\kappa(T)=\kappa_0e^{-\frac{E_{\text{G}}}{2kT}}
		\]

	\section{Bandabstand $E_{\text{g}}$}
		\[
			E_{\text{g}}=-2k\frac{\frac{\kappa_1}{\kappa_2}}{\frac{1}{T_1}-\frac{1}{T_2}}=\frac{hc}{\lambda_{\mathrm{Grenz}}}=h\cdot f_{\mathrm{Grenz}}
		\]
		\begin{table}[h]
		\begin{tabular}{ll}
		$hc=h'=1,24\mu$meV \\
		\end{tabular}
		\end{table}

\chapter{Wichtige Konstanten}
	\begin{table}[ht]
	\begin{tabularx}{\textwidth}{lll}
	Abkürzung & Konstante & Konstantenname\\
	\toprule
	$h$ & $6,626\cdot10^{-34}\mathrm{Js}=4,136\cdot10^{-15}\mathrm{eVs}$ & Plancksches Wirkungsquantum\\
	\midrule
	$h'$ & $hc=1,24\mu\mathrm{meV}$ & \\
	\midrule
	$\bar{h}$ & $1,055\cdot10^{-34}\mathrm{Js}=6,58\cdot10^{-16}\mathrm{eVs}$ & Drehimpulsquantum\\
	\midrule
	$m_0$ & $9,109\cdot10^{-31}$kg & Ruhemasse Elektron\\
	\midrule
	$k$ & $1,381\cdot10^{-23}\nicefrac{\mathrm{J}}{\mathrm{K}}=8,617\cdot10^{-5}\nicefrac{\mathrm{eV}}{\mathrm{K}}$& Boltzmann-Konstante\\
	\midrule
	$e$ & $1,602\cdot10^{-19}$C [C=As] & Elementarladung\\
	\midrule
	1eV & $e=1,602\cdot10^{-19}$J & \\
	\midrule
	$\varepsilon_0$ & $8,854\cdot10^{-12}\nicefrac{\mathrm{F}}{\mathrm{m}}=\nicefrac{\mathrm{As}}{\mathrm{Vm}}$ & Dielektrizitäts-/Permittivitätskonst.\\
	\midrule
	$\mu_0$ & $1,256\cdot10^{-6}\nicefrac{\mathrm{H}}{\mathrm{m}} \left[\nicefrac{\mathrm{Vs}}{\mathrm{Am}}\right]$&Magnetische Feldkonstante\\
	\midrule
	$c_0$ & $2,9979\cdot10^8\nicefrac{\mathrm{m}}{\mathrm{s}}$ & Lichtgeschwindigkeit\\
	\bottomrule
	\end{tabularx}
	\label{Konstanten}
	\caption{Wichtige Konstanten, Quelle: HMS, 12. Auflage}
	\end{table}

\chapter{Einheiten}
	\begin{table}[ht]
	\begin{tabularx}{\textwidth}{lll}
	Größe & Abkürzung & Einheiten\\
	\toprule
	Elektrische Feldkonstante & $\varepsilon_0$ 
	& $\frac{\mathrm{As}}{\mathrm{m}}=\frac{\mathrm{C}}{\mathrm{Vm}}=\frac{\mathrm{C}^2}{\mathrm{Nm}^2}$\\
	\midrule
	Elektrische Feldstärke & $E$ 
	& $\frac{\mathrm{V}}{\mathrm{m}}=\frac{\mathrm{VN}}{\mathrm{Ws}}=\frac{\mathrm{N}}{\mathrm{As}}=\frac{\mathrm{N}}{\mathrm{C}}=\frac{\mathrm{kgm}^2}{\mathrm{s}^3\mathrm{A}}$\\
	\midrule
	Elektrische Kapazität & $C$ 
	& $F=\frac{\mathrm{C}}{\mathrm{V}}=\frac{\mathrm{As}}{\mathrm{V}}=\frac{\mathrm{s}}{\Omega}=\frac{\mathrm{As}^4}{\mathrm{kgm}^2}=\frac{\mathrm{Nm}}{\mathrm{V}^2}=\frac{\mathrm{kgm}^2}{\mathrm{V}^2\mathrm{s}^2}$\\
	\midrule
	Elektrische Ladung & $Q$ & C = AS = $\frac{\mathrm{Ws}}{\mathrm{V}}=\frac{\mathrm{Nm}}{\mathrm{V}}=\frac{\mathrm{kgm}^2}{\mathrm{Vs}^2}$\\
	\midrule
	Elektrische Spannung & $U$ & V = A$\Omega = \frac{\mathrm{W}}{\mathrm{A}} = \frac{\mathrm{Nm}}{\mathrm{As}} = \frac{\mathrm{Ws}}{\mathrm{As}} = \frac{\mathrm{J}}{\mathrm{As}} = \frac{\mathrm{kgm}^2}{\mathrm{as}^3}$\\
	\midrule
	Elektrische Stromstärke & $I$ & A = $\frac{\mathrm{V}}{\Omega}$\\
	\midrule
	Elektrischer Widerstand & $R$ & $\Omega=\frac{\mathrm{V}}{\mathrm{A}}=\frac{\mathrm{kgm}^2}{\mathrm{A}^2\mathrm{s}^3}$\\
	\midrule
	Energie & $E$/$W$ & J = VAs = Ws = Nm = $\frac{\mathrm{kgm}^2}{\mathrm{s}^2}=\frac{1}{1,602\cdot10^{-19}}$eV\\
	\midrule
	Kraft & $F$ & N = $\frac{\mathrm{VAs}}{\mathrm{m}}=\frac{\mathrm{Ws}}{\mathrm{m}}=\frac{\mathrm{J}}{\mathrm{m}}=\frac{\mathrm{kgm}}{\mathrm{s}^2}$\\
	\midrule
	Leistung & $P$ & W = VA = $\frac{\mathrm{J}}{\mathrm{s}}=\frac{\mathrm{Nm}}{\mathrm{s}}=\frac{\mathrm{kgm}^2}{\mathrm{s}^3}$\\
	\midrule
	Magnetische Flussdichte & $B$ & T = $\frac{\mathrm{kg}}{\mathrm{As}^2}=\frac{\mathrm{Vs}}{\mathrm{m}^2}$\\
	\bottomrule
	\end{tabularx}
	\label{Einheiten}
	\caption{Einheitentabelle, Quelle: HMS, 12. Auflage}
	\end{table}

\input{content/chapters/hinweise}
\input{content/chapters/diagramme}
\begin{table}[here]
\begin{tabular}{ll}
\toprule
  $\Phi=\Phi_0\cdot e^{-a\cdot d}$&	Lichtleistung\\
\toprule
  $\Phi \ldots$ Strahlungsleistung/Photonenfluss & z.B. Lichtleistung beim Austritt aus einer Glasfaser \\
  $\Phi_0 \ldots$ reine Strahlungsleistung & also vor dem Einspeisen in die Glasfaser\\
  $d \ldots$ Eindringtiefe& hier ist das die Leitungslänge \\
  $a \ldots$ Absorptionskoeffizient & Dämpfung innerhalb der Glasfaser\\
\midrule
\end{tabular}
\end{table}

\section*{verschiedene Werte nach SI-Einheiten}
\begin{table}[here]
\begin{tabular}{llll}
Wert			 & SI-Einheiten				&Einheit& Art			\\\toprule
$1N_{ewton}		$&$ \frac{kg\cdot m}{s^2}		$&[1J]	&  Kraft		\\\midrule
$1Ns			$&$ \frac{kg\cdot m}{s}			$&	& 	 		\\\midrule
$1Nm			$&$ \frac{kg\cdot m^2}{s^2}		$&	&  Arbeit 		\\\midrule
$1Nms			$&$ \frac{kg\cdot m^2}{s}		$&	&  Drehimpuls		\\\midrule
$?$			 &$ kg\cdot m^2				$&	&  Massenträgheitsmoment\\\midrule
$1\frac{N}{m}		$&$ \frac{kg}{s^2}			$&	&  			\\\midrule
$1\frac{N}{m^2}		$&$ \frac{kg}{m\cdot s^2}		$&	&  Druck		\\\midrule
$1\frac{m}{V}		$&$ \frac{kg}{m^3}			$&	&  Dicht		\\\midrule
$1W_{att}		$&$ \frac{kg\cdot m}{s^2}		$&	&  Leistung		\\\midrule
$1Ws			$&$ \frac{kg\cdot m^2}{s^2}		$&	&  Energie		\\\midrule
$1V_{olt}		$&$ \frac{kg\cdot m^2}{A\cdot s^3}	$&	&  SpannungWiderstand	\\\midrule
$1Vs			$&$ \frac{kg\cdot m^2}{A\cdot s^2}	$&	&  			\\\midrule
$1\frac{V}{m}		$&$ \frac{kg\cdot m}{A\cdot s^3}	$&	&  Feldstärke		\\\midrule
$1\Omega		$&$ \frac{kg\cdot m}{s^2}		$&	&  Widerstand		\\\midrule
$1A_{mpere}		$&$ \frac{kg\cdot m}{s^2}		$&	&  Stromstärke		\\\midrule
$1F_{arad}		$&$ \frac{kg\cdot m}{s^2}		$&	& 			\\\midrule
$1H_{enry}		$&$ \frac{kg\cdot m^2}{A^2\cdot s^2}	$&	& 			\\\midrule
$1C_{ouloumb}		$&$ \frac{kg\cdot m}{s^2}		$&	& 			\\\midrule
$1eV			$&$ \frac{A^2\cdot s^4}{kg\cdot m^2}	$&	& 			\\\midrule
$1B			$&$ \frac{kg}{A\cdot s^2}		$&	&  Magnetfeldstärke	\\\midrule
\end{tabular}
\end{table}


% Hier kommmen Tabellen und Bilder
\listoftables
%\listoffigures

\end{document}
%% Dokument ENDE %%%%%%%%%%%%%%%%%%%%%%%%%%%%%%%%%%%%%%%%%%%%%%%%%%%%%%%%%%

